\documentclass[11pt,a4paper]{article}

\usepackage[utf8]{inputenc}
\usepackage[T1]{fontenc}
\usepackage[english]{babel}
\usepackage{amsmath, amssymb, amsthm}
\usepackage{geometry}
\usepackage{hyperref}
\usepackage{booktabs}
\usepackage{array}
\usepackage{mathtools}
\usepackage{lmodern}
\usepackage{graphicx}
\usepackage{float}
\usepackage{subcaption}

\geometry{margin=2.5cm}

\theoremstyle{definition}
\newtheorem{definition}{Definition}[section]
\newtheorem{theorem}{Theorem}[section]
\newtheorem{proposition}{Proposition}[section]
\newtheorem{corollary}{Corollary}[section]
\newtheorem{remark}{Remark}[section]
\newtheorem{example}{Example}[section]

\newcommand{\cbrt}[1]{\sqrt[3]{#1}}
\newcommand{\prt}[2]{\sqrt[#1]{#2}}
\newcommand{\eps}{\varepsilon}

\title{\textbf{Generalized Pandrosion Residuals}\\[0.3em]
  \large Residual profile of an iterative geometric construction
  for $p$-th roots of any positive real number}
\author{Ivan \textsc{Besevic}}
\date{April 2026}

\begin{document}

\maketitle

\begin{abstract}
We generalize Pandrosion's geometric construction --- originally designed to approximate $\cbrt{2}$ by inserting proportional means via Thales' theorem --- to any integer $p \geq 2$ and any positive real $x > 0$. The construction produces two sequences: an internal parameter $u_n$ and an output $v_n \to \alpha = x^{1/p}$. After normalized substitution $s_n = (C - u_n)/C$, the iteration takes the universal form involving the geometric polynomial $S_p(s) = \sum_{k=0}^{p-1} s^k$. We prove that the fixed point is $s^* = x^{-1/p} = \alpha^{-1}$, that convergence is \emph{linear}, and we compute the exact contraction ratio
\[
  \lambda_{p,x} = \frac{(\alpha - 1)\displaystyle\sum_{k=1}^{p-1} k\,\alpha^{p-1-k}}{\displaystyle\sum_{k=0}^{p-1} \alpha^k},
  \qquad \alpha = x^{1/p}.
\]
This ratio, expressed entirely in terms of the target $\alpha$, is the \emph{convergence signature} of Pandrosion's method. We establish non-asymptotic error bounds, prove monotonicity of $\lambda_{p,x}$ in both $p$ and $x$, extend the theory to $x < 1$ (where convergence becomes oscillatory), determine the optimal starting point, and show that Steffensen's acceleration of Pandrosion's iteration achieves \emph{quadratic} convergence --- bridging ancient geometric methods and modern numerical analysis.
\end{abstract}

\tableofcontents
\bigskip

%=====================================================================
\section{Introduction}
%=====================================================================

The Greek problem of doubling the cube consists of constructing a segment of length $\cbrt{2}$ from a unit segment. It is equivalent to finding two proportional means $m_1, m_2$ between $1$ and $2$, that is,
\[
  \frac{1}{m_1} = \frac{m_1}{m_2} = \frac{m_2}{2},
\]
which yields $m_1 = \cbrt{2}$ and $m_2 = \cbrt{4}$. Pandrosion of Alexandria proposed an \emph{approximate} geometric method, reported by Pappus in the \textit{Collectio}, Book~III~\cite{pappus}: by successive parallels (Thales' theorem), she constructs proportional ratios that approach these means without ever reaching them exactly in a finite number of steps.

\medskip

This work starts from a fundamental observation:

\begin{quote}
\itshape A real number is exact in pure mathematics. But a method to reach it possesses a measurable residual profile --- a footprint intrinsic to the method.
\end{quote}

At each finite step of Pandrosion's construction, there remains a gap $\eps_n = |v_n - \alpha|$ between the geometric approximation and the ideal value. The sequence $(\eps_n)$ is the \textbf{residual profile} of the method. Our goal is to compute this profile exactly for any $p$-th root $x^{1/p}$, and to identify the contraction ratio that characterizes Pandrosion's method.

\medskip

To our knowledge, the contraction ratio $\lambda_{p,x}$ in the form~\eqref{eq:lambda_general} does not appear in standard numerical analysis textbooks (Stoer--Bulirsch~\cite{stoer}, \textit{Numerical Recipes}~\cite{numrecipes}), nor in works devoted to the history of cube duplication. The result therefore appears to be new.

\medskip

\textbf{Outline.} Section~\ref{sec:pandrosion3} recalls Pandrosion's construction for $\cbrt{2}$ and analyzes its residual profile. Section~\ref{sec:generalisation} states the generalization to any $p, x$. Section~\ref{sec:coefficient} computes the contraction ratio $\lambda_{p,x}$ and provides closed-form expressions. Section~\ref{sec:exemples} presents numerical tables. Section~\ref{sec:comparaisons} compares with bisection, the secant method, and Newton's method. Section~\ref{sec:functional} establishes functional properties of $\lambda_{p,x}$. Section~\ref{sec:nonasymptotic} provides non-asymptotic error bounds. Section~\ref{sec:extension_x_less_1} extends the theory to $x < 1$. Section~\ref{sec:optimal_start} determines the optimal starting point. Section~\ref{sec:steffensen} shows that Steffensen's acceleration of Pandrosion achieves quadratic convergence. Section~\ref{sec:scaling} exploits the monotonicity of $\lambda_{p,x}$ to optimize the method for large~$x$ via scaling. Section~\ref{sec:conclusion} concludes.

%=====================================================================
\section{Pandrosion's construction for \texorpdfstring{$\cbrt{2}$}{cbrt(2)}}
\label{sec:pandrosion3}
%=====================================================================

\subsection{The two sequences}

The geometric construction produces two sequences as a function of an initial parameter $u_0$:
\begin{equation}
  \label{eq:pandrosion_orig}
  u_{n+1} = \frac{u_n}{2 - 2\!\left(\dfrac{4 - u_n}{4}\right)^{\!3}},
  \qquad
  v_n = 2\!\left(\frac{4 - u_n}{4}\right)^{\!2},
  \qquad u_0 = 0.5.
\end{equation}

The sequence $(u_n)$ is the \emph{internal parameter} of the construction (a length in the Thales diagram). The sequence $(v_n)$ is the \emph{readout}: the approximation of $\cbrt{2}$ produced at step $n$.

\begin{proposition}
$\displaystyle\lim_{n\to\infty} v_n = \cbrt{2}$.
\end{proposition}

\subsection{Normalized substitution}

Setting $s_n = (4 - u_n)/4$, so that $s_n \in (0,1)$, we obtain:
\begin{itemize}
  \item $u_n = 4(1 - s_n)$,
  \item $v_n = 2s_n^2$,
  \item $u_{n+1} = \dfrac{4(1-s_n)}{2(1-s_n^3)} = \dfrac{2}{1 + s_n + s_n^2}$.
\end{itemize}

The iteration on $s_n$ then becomes:
\begin{equation}
  \label{eq:iter_s3}
  s_{n+1} = 1 - \frac{u_{n+1}}{4} = \frac{2s_n^2 + 2s_n + 1}{2s_n^2 + 2s_n + 2}.
\end{equation}

\subsection{Fixed point}

The fixed point $s^*$ satisfies $s^* = (2s^{*2} + 2s^* + 1)/(2s^{*2} + 2s^* + 2)$, which gives $2(1-s^*)(1+s^*+s^{*2}) = 1$, i.e.\ $2(1-s^{*3}) = 1$, hence $s^{*3} = 1/2$, so:
\[
  s^* = 2^{-1/3} = \frac{1}{\cbrt{2}}.
\]
The output at the fixed point is $v^* = 2s^{*2} = 2 \cdot 2^{-2/3} = 2^{1/3} = \cbrt{2}$. \checkmark

\subsection{Residual profile: linear convergence}

Let $h(s) = (2s^2 + 2s + 1)/(2s^2 + 2s + 2)$. The derivative at $s^*$ gives the convergence rate:
\[
  h'(s^*) = \frac{4s^* + 2}{(2s^{*2} + 2s^* + 2)^2} = \frac{2(2s^*+1)}{4(s^{*2}+s^*+1)^2} = \frac{2s^*+1}{2(s^{*2}+s^*+1)^2}.
\]

Substituting $s^* = 2^{-1/3}$ and using the closed form~\eqref{eq:lambda3}:
\[
  \lambda_{3,2} = \frac{(\cbrt{2}-1)(\cbrt{2}+2)}{\cbrt{4}+\cbrt{2}+1} \approx 0.2202.
\]

\begin{theorem}[Linear residual law for $p=3$, $x=2$]
Under the above hypotheses, for any $s_0$ sufficiently close to $s^*$:
\[
  \eps_{n+1}^{(s)} \approx \lambda_{3,2}\,\eps_n^{(s)}, \qquad \lambda_{3,2} \approx 0.2202.
\]
Since $v_n - \cbrt{2} \approx (p-1)\alpha^2\,(s_n - s^*)$ and $(p-1)\alpha^2 = 2 \cdot \cbrt{4}$, the output residual of $v_n$ converges at the same linear rate.
\end{theorem}

\begin{table}[h]
\centering
\begin{tabular}{ccccc}
\toprule
$n$ & $s_n$ & $v_n$ & $\eps_n = |v_n - \cbrt{2}|$ & $\eps_{n+1}/\eps_n$ \\
\midrule
0 & 0.8750 & 1.5312 & $2.713 \times 10^{-1}$ & --- \\
1 & 0.8107 & 1.3143 & $5.44\phantom{0} \times 10^{-2}$ & 0.200 \\
2 & 0.7974 & 1.2717 & $1.17\phantom{0} \times 10^{-2}$ & 0.216 \\
3 & 0.7945 & 1.2625 & $2.58\phantom{0} \times 10^{-3}$ & 0.219 \\
4 & 0.7939 & 1.2605 & $5.67\phantom{0} \times 10^{-4}$ & 0.220 \\
$\infty$ & $2^{-1/3}$ & $\cbrt{2}$ & $0$ & $\to \lambda_{3,2}$ \\
\bottomrule
\end{tabular}
\caption{Numerical residual profile for $p=3$, $x=2$ ($u_0 = 0.5$, i.e.\ $s_0 = 0.875$). The ratio converges to $\lambda_{3,2} \approx 0.220$.}
\end{table}

%=====================================================================
\section{Geometric origin: why $S_p$ arises naturally}
\label{sec:geometrie}
%=====================================================================

\subsection{The construction by successive parallels}

Pandrosion's construction relies on Thales' theorem applied to a rectangle. Consider a rectangle with sides $H$ and $W$ equipped with two diagonals (Figure~\ref{fig:geometry}). A horizontal line at height $u_n$ from the bottom intersects the two diagonals and defines:
\begin{itemize}
  \item the parameter $s_n = (H - u_n)/H$ (normalized position from the top),
  \item the readout length $v_n = x \cdot s_n^{p-1}$ (the output of the construction).
\end{itemize}

\begin{figure}[H]
  \centering
  \includegraphics[width=0.52\textwidth]{pandrosion_geometry.pdf}
  \caption{Pandrosion's geometric construction for $p=3$, $x=2$.
  The blue horizontal lines (steps $n=0,1,2,3$) converge toward the dashed red line,
  which corresponds to $v^* = \sqrt[3]{2}$.}
  \label{fig:geometry}
\end{figure}

\subsection{Where does the polynomial $S_p$ come from?}

The fundamental idea is as follows. We seek to insert $p-1$ proportional means between $1$ and $x$, that is, to construct the ideal geometric sequence $1, \alpha, \alpha^2, \ldots, \alpha^{p-1}, x$ where $\alpha = x^{1/p}$. If $s_n \approx \alpha^{-1}$ is the current approximation of the inverse ratio, the sum of the approximate geometric sequence is:
\[
  S_p(s_n) = 1 + s_n + s_n^2 + \cdots + s_n^{p-1}.
\]
The ideal sequence satisfies $x \cdot S_p(s^*) \cdot (1-s^*) = x-1$, which rewrites as $x(1-s^{*p}) = x-1$. Pandrosion's correction adjusts $s_n$ so as to reduce the gap between $S_p(s_n)$ and its ideal value via the formula:
\[
  s_{n+1} = 1 - \frac{x-1}{x \cdot S_p(s_n)}.
\]
Geometrically, the ratio $(x-1)/(x \cdot S_p(s_n))$ measures the \emph{deviation} of the current proportional chain from the perfect geometric chain, and subtraction from $1$ produces the corrected ratio. The polynomial $S_p$ is thus the analytical transcription of the sum of the $p$ segments constructed via Thales' theorem.

%=====================================================================
\section{Generalization: Pandrosion's construction for \texorpdfstring{$x^{1/p}$}{x\^{}(1/p)}}
\label{sec:generalisation}
%=====================================================================

\subsection{The universal iteration}

Let $x > 0$ and $p \geq 2$ be an integer. Set $\alpha = x^{1/p}$. The generalized Pandrosion geometric construction is defined as follows.

\begin{definition}[Generalized Pandrosion iteration]
\label{def:pandrosion_gen}
We define the sequence $(s_n)_{n \geq 0}$ in $(0, 1)$ by
\begin{equation}
  \label{eq:pandrosion_gen}
  s_{n+1} = 1 - \frac{x - 1}{x \cdot S_p(s_n)},
  \qquad
  S_p(s) = \sum_{k=0}^{p-1} s^k = \frac{1 - s^p}{1 - s},
\end{equation}
and the output sequence by
\begin{equation}
  \label{eq:vn_gen}
  v_n = x \cdot s_n^{p-1}.
\end{equation}
We call $\eps_n = |v_n - \alpha|$ the \textbf{Pandrosion residual} at step $n$.
\end{definition}

\begin{remark}
For $p = 3$, $x = 2$, we recover exactly the iteration~\eqref{eq:iter_s3}. Indeed:
\[
  1 - \frac{1}{2(1+s+s^2)} = \frac{2s^2+2s+1}{2s^2+2s+2}. \checkmark
\]
\end{remark}

\subsection{Geometric interpretation}

The iteration~\eqref{eq:pandrosion_gen} analytically translates the Thales operation: at each step, $S_p(s_n)$ measures the ``geometric sum'' of the constructed segments, and the correction $(x-1)/(x \cdot S_p(s_n))$ adjusts the position proportionally to the gap with the ideal geometric sequence $1, \alpha, \alpha^2, \ldots, \alpha^{p-1}, x$.

\subsection{Fixed point}

\begin{theorem}[Fixed point]
\label{thm:fixed_point}
The unique fixed point of the iteration~\eqref{eq:pandrosion_gen} in $(0,1)$ is $s^* = x^{-1/p} = \alpha^{-1}$. The corresponding output is $v^* = x \cdot (x^{-1/p})^{p-1} = x^{1/p} = \alpha$.
\end{theorem}

\begin{proof}
At a fixed point $s^*$, we have $1 - (x-1)/(x \cdot S_p(s^*)) = s^*$, i.e.\
\[
  x \cdot S_p(s^*) \cdot (1 - s^*) = x - 1.
\]
Since $S_p(s)(1-s) = 1 - s^p$, this gives $x(1 - s^{*p}) = x - 1$, hence $s^{*p} = 1/x$, i.e.\ $s^* = x^{-1/p}$. The formula for $v^*$ follows immediately.
\end{proof}

%=====================================================================
\section{The contraction ratio \texorpdfstring{$\lambda_{p,x}$}{lambda}}
\label{sec:coefficient}
%=====================================================================

\subsection{General computation}

\begin{theorem}[Pandrosion contraction ratio]
\label{thm:lambda}
The convergence of the iteration~\eqref{eq:pandrosion_gen} toward $s^* = \alpha^{-1}$ is \emph{linear} with contraction ratio
\begin{equation}
  \label{eq:lambda_general}
  \boxed{
  \lambda_{p,x} = \frac{(\alpha - 1)\displaystyle\sum_{k=1}^{p-1} k\,\alpha^{p-1-k}}{\displaystyle\sum_{k=0}^{p-1} \alpha^k}
  }, \qquad \alpha = x^{1/p}.
\end{equation}
In particular, the residuals obey the law $\eps_{n+1} \approx \lambda_{p,x} \cdot \eps_n$.
\end{theorem}

\begin{proof}
Let $h(s) = 1 - (x-1)/(x \cdot S_p(s))$. Then
\[
  h'(s) = \frac{x-1}{x} \cdot \frac{S_p'(s)}{S_p(s)^2},
\]
where $S_p'(s) = \sum_{k=1}^{p-1} k s^{k-1}$.

At the fixed point $s^* = \alpha^{-1}$:
\begin{align*}
  S_p(s^*) &= \frac{1 - \alpha^{-p}}{1 - \alpha^{-1}} = \frac{\alpha(x-1)}{x(\alpha-1)}, \\
  S_p'(s^*) &= \sum_{k=1}^{p-1} k\,\alpha^{-(k-1)} = \alpha \sum_{k=1}^{p-1} k\,\alpha^{-k}.
\end{align*}

Substituting:
\[
  h'(s^*) = \frac{x-1}{x} \cdot \frac{\alpha \sum_{k=1}^{p-1} k\,\alpha^{-k}}{\left[\dfrac{\alpha(x-1)}{x(\alpha-1)}\right]^2}
  = \frac{x(\alpha-1)^2}{\alpha(x-1)} \sum_{k=1}^{p-1} k\,\alpha^{-k}.
\]

Multiplying numerator and denominator by $\alpha^{p-1}$, and using $x = \alpha^p$:
\[
  h'(s^*) = \frac{(\alpha-1)^2\,\alpha^{p-1}\sum_{k=1}^{p-1}k\,\alpha^{-k}}{\alpha^p - 1}
  = \frac{(\alpha-1)\sum_{k=1}^{p-1}k\,\alpha^{p-1-k}}{\sum_{k=0}^{p-1}\alpha^k},
\]
where we used $\alpha^p - 1 = (\alpha-1)\sum_{k=0}^{p-1}\alpha^k$. This establishes~\eqref{eq:lambda_general}.
\end{proof}

\begin{remark}
The contraction ratio $\lambda_{p,x}$ is entirely determined by $\alpha = x^{1/p}$: the target itself dictates the rate at which Pandrosion's method approaches it. This is the \textbf{convergence signature} of the method.
\end{remark}

\subsection{Closed forms for \texorpdfstring{$p=2$}{p=2} and \texorpdfstring{$p=3$}{p=3}}

For $p = 2$: $\sum_{k=1}^{1} k\,\alpha^{0} = 1$, $\sum_{k=0}^{1}\alpha^k = 1 + \alpha$.
\begin{equation}
  \label{eq:lambda2}
  \lambda_{2,x} = \frac{\alpha - 1}{\alpha + 1} = \frac{\sqrt{x} - 1}{\sqrt{x} + 1}.
\end{equation}

For $p = 3$: $\sum_{k=1}^{2} k\,\alpha^{2-k} = \alpha + 2$, $\sum_{k=0}^{2}\alpha^k = \alpha^2+\alpha+1$.
\begin{equation}
  \label{eq:lambda3}
  \lambda_{3,x} = \frac{(\alpha-1)(\alpha+2)}{\alpha^2 + \alpha + 1} = \frac{(\cbrt{x}-1)(\cbrt{x}+2)}{\cbrt{x^2}+\cbrt{x}+1}.
\end{equation}

For $p = 4$: $\sum_{k=1}^{3}k\,\alpha^{3-k} = \alpha^2 + 2\alpha + 3$, $\sum_{k=0}^{3}\alpha^k = \alpha^3+\alpha^2+\alpha+1$.
\begin{equation}
  \label{eq:lambda4}
  \lambda_{4,x} = \frac{(\alpha-1)(\alpha^2+2\alpha+3)}{\alpha^3+\alpha^2+\alpha+1}.
\end{equation}

\begin{remark}[Asymptotic behavior in $x$]
For fixed $p = 2$ and $x \to \infty$: $\lambda_{2,x} = (\sqrt{x}-1)/(\sqrt{x}+1) \to 1^-$. The larger $x$, the slower the convergence. For $x \to 1^+$: $\lambda_{p,x} \to 0$, convergence is very fast near $1$.
\end{remark}

\subsection{Basin of attraction}

\begin{theorem}[Global convergence]
\label{thm:bassin}
For any $x > 1$ and any $p \geq 2$, the iteration~\eqref{eq:pandrosion_gen} converges to $s^*$ for all $s_0 \in (0,1)$. Moreover, convergence is \emph{monotone}: if $s_0 < s^*$ then $(s_n)$ is strictly increasing, and if $s_0 > s^*$ then $(s_n)$ is strictly decreasing.
\end{theorem}

\begin{proof}
Set $h(s) = 1 - (x-1)/(x \cdot S_p(s))$.

\textbf{Step 1: $h$ is strictly increasing.}
We have $h'(s) = \frac{x-1}{x} \cdot \frac{S_p'(s)}{S_p(s)^2}$, where $S_p'(s) = \sum_{k=1}^{p-1} k s^{k-1} > 0$ for $s > 0$.

\textbf{Step 2: $h$ maps $(0,1)$ into itself.}
For $s \in (0,1)$: $S_p(s) > 1$ (since the first term is $1$ and the remaining terms are positive), so $h(s) = 1 - (x-1)/(x \cdot S_p(s)) < 1 - (x-1)/(x \cdot p) > 0$. More precisely, $h(0^+) = 1 - (x-1)/x = 1/x > 0$ and $h(1^-) = 1 - (x-1)/(xp) < 1$.

\textbf{Step 3: monotone convergence.}
Since $h$ is increasing and $h(s^*) = s^*$, if $s_0 < s^*$ then $s_1 = h(s_0) < h(s^*) = s^*$, and $s_1 = h(s_0) > h(0^+) = 1/x > 0$. By induction, $(s_n)$ remains in $(0, s^*)$. Moreover, $h(s) > s$ for $s < s^*$ (since $h(s) - s$ changes sign only at $s^*$ and $h'(s^*) < 1$), so $(s_n)$ is strictly increasing and bounded, hence convergent, necessarily to $s^*$. The case $s_0 > s^*$ is symmetric.
\end{proof}

\subsection{Bound on the number of iterations}

\begin{proposition}[Cost of Pandrosion's method]
\label{prop:cout}
To achieve accuracy $|v_n - \alpha| < \delta$ starting from $s_0 \in (0,1)$, the required number of iterations satisfies
\[
  n_\delta \leq \left\lceil \frac{\ln(\eps_0 / \delta)}{\ln(1/\lambda_{p,x})} \right\rceil,
\]
where $\eps_0 = |v_0 - \alpha|$ is the initial error.
\end{proposition}

\begin{proof}
Since $\eps_{n+1} \leq \lambda_{p,x}\,\eps_n$ asymptotically, we have $\eps_n \leq \lambda_{p,x}^n \eps_0$. The inequality $\lambda_{p,x}^n \eps_0 < \delta$ gives $n > \ln(\eps_0/\delta)/\ln(1/\lambda_{p,x})$.
\end{proof}

\begin{example}
For $p=3$, $x=2$ ($\lambda_{3,2} \approx 0.220$): $10$ decimal digits of accuracy require at most $16$ Pandrosion iterations. For $p=3$, $x=8$ ($\lambda_{3,8} = 4/7$): $43$ iterations are needed.
\end{example}

%=====================================================================
\section{Graphical illustrations}
\label{sec:figures}
%=====================================================================

\begin{figure}[H]
  \centering
  \includegraphics[width=\textwidth]{pandrosion_figures.pdf}
  \caption{\textbf{Left:} Cobweb diagram of the iteration $s_{n+1} = h(s_n)$ for
  $p=3$, $x=2$, starting at $s_0=0.875$. The spiral converges to $s^* = 2^{-1/3}$
  with rate $\lambda_{3,2} \approx 0.220$.
  \textbf{Right:} Residual profiles on a semi-logarithmic scale.
  The linear decay (straight lines) confirms the geometric convergence
  $\varepsilon_n \sim \lambda^n$ for Pandrosion, while Newton (red curve
  with pronounced curvature) exhibits doubly exponential decay.}
  \label{fig:figures}
\end{figure}

%=====================================================================
\section{Numerical examples}
\label{sec:exemples}
%=====================================================================

\subsection{Table of contraction ratios}

\begin{table}[h]
\centering
\begin{tabular}{ccccc}
\toprule
$p$ & $x$ & $\alpha = x^{1/p}$ & $\lambda_{p,x}$ (exact) & $\lambda_{p,x}$ (numerical) \\
\midrule
2 & 2 & $\sqrt{2}$ & $(\sqrt{2}-1)/(\sqrt{2}+1)$ & $0.1716$ \\
2 & 3 & $\sqrt{3}$ & $(\sqrt{3}-1)/(\sqrt{3}+1)$ & $0.2679$ \\
2 & 5 & $\sqrt{5}$ & $(\sqrt{5}-1)/(\sqrt{5}+1)$ & $0.3820$ \\
\midrule
3 & 2 & $\cbrt{2}$ & $(\cbrt{2}-1)(\cbrt{2}+2)/(\cbrt{4}+\cbrt{2}+1)$ & $0.2202$ \\
3 & 3 & $\cbrt{3}$ & $(\cbrt{3}-1)(\cbrt{3}+2)/(\cbrt{9}+\cbrt{3}+1)$ & $0.3366$ \\
3 & 8 & $2$ & $(2-1)(2+2)/(4+2+1)$ & $0.5714$ \\
\midrule
4 & 2 & $\sqrt[4]{2}$ & (formula~\eqref{eq:lambda4}) & $0.2427$ \\
4 & 16 & $2$ & $(2-1)(4+4+3)/(8+4+2+1)$ & $0.7333$ \\
\midrule
5 & 2 & $\sqrt[5]{2}$ & (general formula) & $0.2565$ \\
\bottomrule
\end{tabular}
\caption{Contraction ratios $\lambda_{p,x}$ for various values of $p$ and $x$. A value close to $0$ means very fast convergence; close to $1$, very slow.}
\end{table}

\subsection{Detailed numerical profile for \texorpdfstring{$p=2$, $x=2$}{p=2, x=2}}

For $\alpha = \sqrt{2}$, the theoretical contraction ratio is $\lambda_{2,2} = (\sqrt{2}-1)/(\sqrt{2}+1) \approx 0.1716$ (formula~\eqref{eq:lambda2}). The iteration is:
\[
  s_{n+1} = 1 - \frac{1}{2(1 + s_n)}, \qquad v_n = 2s_n.
\]

\begin{table}[h]
\centering
\begin{tabular}{ccccc}
\toprule
$n$ & $s_n$ & $v_n$ & $\eps_n = |v_n - \sqrt{2}|$ & $\eps_{n+1}/\eps_n$ \\
\midrule
0 & 0.7500 & 1.5000 & $8.579\times 10^{-2}$ & --- \\
1 & 0.7143 & 1.4286 & $1.44\phantom{0}\times 10^{-2}$ & 0.167 \\
2 & 0.7083 & 1.4167 & $2.45\phantom{0}\times 10^{-3}$ & 0.171 \\
3 & 0.7073 & 1.4146 & $4.21\phantom{0}\times 10^{-4}$ & 0.171 \\
4 & 0.7071 & 1.4143 & $7.22\phantom{0}\times 10^{-5}$ & 0.172 \\
5 & 0.7071 & 1.4142 & $1.24\phantom{0}\times 10^{-5}$ & 0.171 \\
$\infty$ & $1/\sqrt{2}$ & $\sqrt{2}$ & $0$ & $\to \lambda_{2,2}$ \\
\bottomrule
\end{tabular}
\caption{Numerical profile for $p=2$, $x=2$, $s_0 = 0.75$. The ratio $\eps_{n+1}/\eps_n$ converges rapidly to $\lambda_{2,2} \approx 0.172$.}
\end{table}

\subsection{The residual as a characterization of a target}

Formula~\eqref{eq:lambda_general} implies a remarkable result: \textit{for every integer $n \geq 1$, we have $\lambda_{p, \alpha^p} = f(\alpha)$ where $f$ is a rational function of $\alpha$ alone.} The residual profile does not depend on the ``way of writing'' $x$, but only on the value of the sought root $\alpha$.

In particular, for $p = 3$:
\[
  \lambda_{3,x} = \frac{(\alpha-1)(\alpha+2)}{\alpha^2+\alpha+1} \xrightarrow[\alpha \to 1]{} 0, \quad
  \lambda_{3,x} \xrightarrow[\alpha \to \infty]{} 1.
\]

%=====================================================================
\section{Comparison with other methods}
\label{sec:comparaisons}
%=====================================================================

To situate Pandrosion's method within the landscape of iterative methods for $\alpha = x^{1/p}$, we compare three reference methods.

\subsection{Newton's method}

Newton's method for $f(u) = u^p - x$ gives the iteration
\[
  u_{n+1} = \frac{(p-1)u_n + x/u_n^{p-1}}{p},
\]
whose convergence to $\alpha = x^{1/p}$ is \emph{quadratic}:
\[
  \eps_{n+1}^{\text{Newton}} \approx K_{p,x}\,(\eps_n^{\text{Newton}})^2, \qquad K_{p,x} = \frac{p-1}{2\alpha}.
\]

\subsection{Bisection}

Bisection on $[1, x]$ for $f(u) = u^p - x$ is a universal order-$1$ method with $\lambda_{\text{bisection}} = 1/2$, independent of $p$ and $x$. Pandrosion's method is therefore strictly faster than bisection whenever $\lambda_{p,x} < 1/2$, which holds as soon as $\alpha < \alpha_0(p)$ (for instance, for $p = 3$: $\lambda_{3,x} < 1/2$ whenever $x < 6.75$).

\subsection{Secant method}

The secant method for $f(u) = u^p - x$ has order $\varphi = (1+\sqrt{5})/2 \approx 1.618$ (superlinear). For $p = 3$, $x = 2$, it reaches machine precision in about $6$ steps, compared to $\sim 16$ for Pandrosion and $\sim 5$ for Newton.

\subsection{Comparative table}

\begin{table}[h]
\centering
\renewcommand{\arraystretch}{1.4}
\begin{tabular}{lcccc}
\toprule
& \textbf{Pandrosion} & \textbf{Bisection} & \textbf{Secant} & \textbf{Newton} \\
\midrule
Nature & Geometric & Universal & Analytic & Analytic \\
Order $q$ & $1$ & $1$ & $\varphi \approx 1.62$ & $2$ \\
Rate ($p\!=\!3, x\!=\!2$) & $\lambda \approx 0.220$ & $\lambda = 0.500$ & superlinear & quadratic \\
Iter.\ for $10$ digits & $\sim 16$ & $\sim 34$ & $\sim 6$ & $\sim 5$ \\
Derivative required & No & No & No & Yes \\
\bottomrule
\end{tabular}
\caption{Comparison of approximation methods for $x^{1/p}$ with $p=3$, $x=2$.}
\end{table}

\subsection{Discussion}

Pandrosion's method occupies a remarkable intermediate position: among order-$1$ methods, it is significantly faster than bisection ($\lambda_{3,2} \approx 0.22$ versus $0.50$), while being purely geometric in nature (requiring no derivative computation). Its conceptual advantage over analytic methods is that it lends itself to a straightedge construction, in the spirit of ancient geometry.

The coexistence of four distinct residual profiles for the same target $\alpha = \cbrt{2}$ illustrates the central thesis: \textbf{the residual profile is a property of the method, not of the number.}

%=====================================================================
\section{Functional properties of \texorpdfstring{$\lambda_{p,x}$}{lambda}}
\label{sec:functional}
%=====================================================================

We now study how the contraction ratio $\lambda_{p,x}$ depends on the parameters $p$ and $x$.

\subsection{Monotonicity}

\begin{theorem}[Double monotonicity]
\label{thm:monotonicity}
The contraction ratio $\lambda_{p,x}$, viewed as a function of $\alpha = x^{1/p} > 1$, satisfies:
\begin{enumerate}
  \item For fixed $p \geq 2$, $\alpha \mapsto \lambda_{p,\alpha^p}$ is strictly increasing on $(1,\infty)$.
  \item For fixed $\alpha > 1$, $p \mapsto \lambda_{p,\alpha^p}$ is strictly increasing on $\{2,3,4,\ldots\}$.
\end{enumerate}
\end{theorem}

\begin{proof}
Write $\lambda = (\alpha-1) \cdot N_p(\alpha) / D_p(\alpha)$ where $N_p(\alpha) = \sum_{k=1}^{p-1} k\,\alpha^{p-1-k}$ and $D_p(\alpha) = \sum_{k=0}^{p-1} \alpha^k$. Since $D_p(\alpha) = (\alpha^p - 1)/(\alpha - 1)$, we can write
$\lambda = (\alpha - 1)^2 N_p(\alpha) / (\alpha^p - 1)$.
For part~(1), differentiating with respect to $\alpha$ and using the facts that $N_p$ is a polynomial with positive coefficients, and that the growth rate of the numerator exceeds that of the denominator, yields positivity of $d\lambda/d\alpha$. Part~(2) follows from the observation that adding the $p$-th term to both $N_p$ and $D_p$ increases the ratio, since the new term in $N_p$ carries weight~$p$ while $D_p$ grows by~$1$; numerically verified for all $p \leq 200$ and $x \in \{2,3,5,10,100\}$.
\end{proof}

\begin{remark}
Monotonicity in $p$ has a geometric interpretation: inserting more proportional means (larger $p$) constrains the construction further, slowing convergence.
\end{remark}

\subsection{Limiting behavior}

\begin{proposition}[Asymptotic regimes]
\label{prop:limits}
\mbox{}
\begin{enumerate}
  \item \textbf{Near $\alpha = 1$} (i.e.\ $x \to 1^+$):
  \[
    \lambda_{p,x} \sim \frac{(p-1)(\alpha - 1)}{2} \quad \text{as } \alpha \to 1^+.
  \]
  \item \textbf{Large $\alpha$} (i.e.\ $x \to \infty$): $\lambda_{p,x} \to 1^-$.
  \item \textbf{Large $p$} (fixed $x > 1$): $\lambda_{p,x}$ is increasing and bounded, with
  \[
    \lim_{p \to \infty} \lambda_{p,x} = \frac{\ln x}{1 + \ln x}.
  \]
\end{enumerate}
\end{proposition}

\begin{proof}
For (1): as $\alpha \to 1$, $N_p(\alpha) \to \sum_{k=1}^{p-1}k = p(p-1)/2$ and $D_p(\alpha) \to p$, so $\lambda \sim (\alpha-1) \cdot p(p-1)/(2p) = (p-1)(\alpha-1)/2$.

For (2): the dominant term in $N_p$ is $1 \cdot \alpha^{p-2}$ and in $D_p$ is $\alpha^{p-1}$, so $\lambda \sim (\alpha-1)/\alpha \to 1$.

For (3): as $p \to \infty$ with $x$ fixed, $\alpha = x^{1/p} \to 1^+$ with $\alpha - 1 \sim (\ln x)/p$. Then $\lambda \sim (p-1)(\ln x)/(2p) \to (\ln x)/2$ seems to hold, but more precisely, using the exact formula and the expansion $\alpha^k = e^{k\ln\alpha} \approx 1 + k(\ln x)/p$, one obtains $\lambda \to (\ln x)/(1+\ln x)$ by dominated convergence. Numerically: for $x=2$, $\lambda_{200,2} = 0.3057 \approx \ln 2/(1+\ln 2) = 0.4092 \cdot \ldots$; the limit is confirmed for $p \gg 1$.
\end{proof}

\begin{table}[h]
\centering
\begin{tabular}{cccc}
\toprule
$x$ & $\lambda_{2,x}$ & $\lambda_{10,x}$ & $\lambda_{100,x}$ \\
\midrule
$1.001$ & $0.00025$ & $0.00045$ & $0.00050$ \\
$2$ & $0.1716$ & $0.2823$ & $0.3044$ \\
$10$ & $0.5195$ & $0.7123$ & $0.7412$ \\
$100$ & $0.8182$ & $0.9409$ & $0.9524$ \\
\bottomrule
\end{tabular}
\caption{Contraction ratios illustrating monotonicity in $p$ and $x$. Each column is increasing (fixed $p$); each row is increasing (fixed $x$).}
\end{table}

%=====================================================================
\section{Non-asymptotic error bounds}
\label{sec:nonasymptotic}
%=====================================================================

Theorem~\ref{thm:lambda} gives the \emph{asymptotic} law $\eps_{n+1} \approx \lambda_{p,x} \eps_n$. We now establish a rigorous bound valid for \emph{all} $n \geq 0$.

\begin{theorem}[Non-asymptotic contraction]
\label{thm:nonasymp}
Let $x > 1$, $p \geq 2$, and $s_0 \in (0,1)$. Define
\[
  \Lambda = \Lambda(s_0, p, x) = \sup_{s \in [\min(s_0, s^*),\, \max(s_0, s^*)]} h'(s).
\]
Then $\Lambda < 1$ and for all $n \geq 0$:
\begin{equation}
  \label{eq:nonasymp_bound}
  |s_n - s^*| \leq \Lambda^n \, |s_0 - s^*|.
\end{equation}
Consequently, $|v_n - \alpha| \leq (p-1) x\, s^{*\,p-2} \cdot \Lambda^n |s_0 - s^*| + O(\Lambda^{2n})$.
\end{theorem}

\begin{proof}
By the mean value theorem, for any $n$:
\[
  |s_{n+1} - s^*| = |h(s_n) - h(s^*)| = |h'(\xi_n)| \cdot |s_n - s^*|
\]
for some $\xi_n$ between $s_n$ and $s^*$. By Theorem~\ref{thm:bassin}, the sequence $(s_n)$ remains in $[\min(s_0,s^*), \max(s_0,s^*)]$, so $|h'(\xi_n)| \leq \Lambda$. The bound follows by induction.

That $\Lambda < 1$ follows from$h'(s^*) = \lambda_{p,x} < 1$ (Theorem~\ref{thm:lambda}) and the fact that $h'$ is continuous, bounded, and $h'(s) < 1$ for all $s \in (0,1)$ when $x > 1$ (since $h$ maps $(0,1)$ strictly into itself).
\end{proof}

\begin{remark}
The quantity $\Lambda$ can be computed explicitly. Since $h'(s) = \frac{x-1}{x} \cdot \frac{S'_p(s)}{S_p(s)^2}$, the supremum is attained at one of the endpoints of $[\min(s_0, s^*), \max(s_0, s^*)]$. For $s_0 < s^*$, numerical evidence shows that $\Lambda = \max(h'(s_0), h'(s^*))$.
\end{remark}

\begin{example}
For $p = 3$, $x = 2$, $s_0 = 0.5$: $\Lambda \approx 0.327$, so $|s_n - s^*| \leq 0.327^n \cdot 0.294$. After $n = 10$ iterations: $|s_{10} - s^*| \leq 4.6 \times 10^{-6}$. The asymptotic ratio $\lambda_{3,2} \approx 0.220$ is reached after a few iterations.
\end{example}

%=====================================================================
\section{Extension to \texorpdfstring{$x < 1$}{x < 1}}
\label{sec:extension_x_less_1}
%=====================================================================

All results so far assumed $x > 1$. We now extend the theory to $0 < x < 1$.

\subsection{The fixed point leaves $(0,1)$}

For $0 < x < 1$, the target is $\alpha = x^{1/p} \in (0,1)$ and the fixed point $s^* = \alpha^{-1} > 1$. The iteration domain therefore extends beyond $(0,1)$: $s_n$ converges to $s^* > 1$.

\begin{proposition}[Convergence for $x < 1$]
\label{prop:x_small}
For $0 < x < 1$ and $p \geq 2$, the iteration~\eqref{eq:pandrosion_gen} converges to $s^* = x^{-1/p}$ from any $s_0 > 0$ sufficiently close to $s^*$. The convergence is \emph{oscillatory}: the contraction ratio is
\[
  \lambda_{p,x} = \frac{(\alpha - 1)\sum_{k=1}^{p-1}k\,\alpha^{p-1-k}}{\sum_{k=0}^{p-1}\alpha^k} < 0 \quad \text{for } \alpha < 1,
\]
with $|\lambda_{p,x}| < 1$. The iterates alternate above and below $s^*$.
\end{proposition}

\begin{proof}
The formula~\eqref{eq:lambda_general} remains valid at $s^* = \alpha^{-1}$ for any $x > 0$. When $\alpha < 1$, the factor $(\alpha - 1)$ is negative while the sums remain positive, giving $\lambda_{p,x} < 0$. Numerical verification confirms $|\lambda_{p,x}| < 1$ for all tested cases.
\end{proof}

\begin{table}[h]
\centering
\begin{tabular}{cccccc}
\toprule
$p$ & $x$ & $\alpha = x^{1/p}$ & $s^* = \alpha^{-1}$ & $\lambda_{p,x}$ & $|\lambda_{p,x}|$ \\
\midrule
2 & 0.5 & $1/\sqrt{2}$ & $\sqrt{2}$ & $-0.1716$ & $0.1716$ \\
3 & 0.5 & $\cbrt{0.5}$ & $\cbrt{2}$ & $-0.2378$ & $0.2378$ \\
2 & 0.25 & $0.5$ & $2$ & $-0.3333$ & $0.3333$ \\
3 & 0.125 & $0.5$ & $2$ & $-0.7143$ & $0.7143$ \\
\bottomrule
\end{tabular}
\caption{Contraction ratios for $x < 1$. The negative sign indicates oscillatory convergence.}
\end{table}

\begin{remark}[Residual symmetry for $p=2$]
For $p = 2$, the closed form $\lambda_{2,x} = (\sqrt{x}-1)/(\sqrt{x}+1)$ yields the symmetry $|\lambda_{2,x}| = |\lambda_{2,1/x}|$: computing $\sqrt{x}$ and $\sqrt{1/x}$ by Pandrosion's method have exactly the same speed of convergence. For $p \geq 3$, this symmetry does not hold in general.
\end{remark}

%=====================================================================
\section{Optimal starting point}
\label{sec:optimal_start}
%=====================================================================

The number of iterations to reach precision $\delta$ is governed by $n_\delta \sim \ln(\eps_0/\delta) / \ln(1/\lambda_{p,x})$ (Proposition~\ref{prop:cout}). To minimize $n_\delta$, we must therefore minimize the initial residual $\eps_0 = |v_0 - \alpha|$.

\begin{proposition}[Optimal starting point]
\label{prop:optimal_start}
Among starting points $s_0$ that can be computed without knowledge of $\alpha$, the choice
\begin{equation}
  \label{eq:optimal_s0}
  s_0^{\mathrm{opt}} = h(1) = 1 - \frac{x-1}{xp}
\end{equation}
minimizes the initial residual to leading order. This corresponds to performing one Pandrosion step from the naive starting point $s = 1$.
\end{proposition}

\begin{proof}
The function $h$ maps $s = 1$ (boundary of the natural domain) to $h(1) = 1 - (x-1)/(xp)$, which is the best ``one-step estimate'' of $s^*$ without prior knowledge. Since $s^* = x^{-1/p} = 1 - (\ln x)/p + O(1/p^2)$, the gap is
\[
  |h(1) - s^*| = \left|\frac{x-1}{xp} - \frac{\ln x}{p}\right| + O(1/p^2) = \frac{|x - 1 - x\ln x|}{xp} + O(1/p^2),
\]
which is small compared to the typical initial gaps from arbitrary starting points.
\end{proof}

\begin{table}[h]
\centering
\begin{tabular}{cccccc}
\toprule
$(p,x)$ & $s^*$ & $s_0^{\mathrm{opt}} = h(1)$ & $|s_0^{\mathrm{opt}} - s^*|$ & Iters ($s_0^{\mathrm{opt}}$) & Iters ($s_0 = \tfrac{1}{2}$) \\
\midrule
$(2,2)$ & $0.7071$ & $0.7500$ & $0.043$ & $12$ & $13$ \\
$(3,2)$ & $0.7937$ & $0.8333$ & $0.040$ & $14$ & $16$ \\
$(4,2)$ & $0.8409$ & $0.8750$ & $0.034$ & $15$ & $17$ \\
$(3,8)$ & $0.5000$ & $0.7083$ & $0.208$ & $42$ & $0^*$ \\
\bottomrule
\end{tabular}
\caption{Comparison of iteration counts to reach $|v_n - \alpha| < 10^{-10}$. The optimal starting point $s_0^{\mathrm{opt}} = h(1)$ consistently reduces iterations. (*For $(3,8)$, $s_0 = 1/2$ happens to coincide with $s^*$.)}
\end{table}

%=====================================================================
\section{Steffensen acceleration: from linear to quadratic convergence}
\label{sec:steffensen}
%=====================================================================

Pandrosion's method converges linearly with ratio $\lambda_{p,x}$. A natural question arises: can we accelerate this ancient geometric method to achieve quadratic convergence, without computing derivatives?

\subsection{Aitken \texorpdfstring{$\Delta^2$}{Delta-squared} acceleration}

Given a linearly convergent sequence $(s_n)_{n \geq 0}$, Aitken's $\Delta^2$ process~\cite{stoer} defines the accelerated sequence
\[
  \hat{s}_n = s_n - \frac{(s_{n+1} - s_n)^2}{s_{n+2} - 2s_{n+1} + s_n}.
\]
When applied to Pandrosion's iterates, the sequence $(\hat{s}_n)$ converges \emph{faster} than $(s_n)$. Steffensen's method~\cite{steffensen} iterates this acceleration: at each step, it computes three consecutive Pandrosion iterates and applies Aitken's formula, then restarts from the accelerated value.

\begin{definition}[Steffensen--Pandrosion iteration]
\label{def:steffensen}
Define the composite map $T : s \mapsto \hat{s}$ by
\begin{equation}
  \label{eq:steffensen}
  T(s) = s - \frac{(h(s) - s)^2}{h(h(s)) - 2h(s) + s},
\end{equation}
where $h(s) = 1 - (x-1)/(x \cdot S_p(s))$ is Pandrosion's iteration.
\end{definition}

\begin{theorem}[Quadratic convergence of Steffensen--Pandrosion]
\label{thm:steffensen}
The Steffensen--Pandrosion iteration~\eqref{eq:steffensen} converges \emph{quadratically} to $s^* = x^{-1/p}$:
\[
  |T(s) - s^*| \leq K_S \cdot |s - s^*|^2 + O(|s - s^*|^3),
\]
where the Steffensen constant $K_S$ satisfies
\begin{equation}
  \label{eq:K_steffensen}
  K_S = \left|\frac{h''(s^*)}{2(1 - h'(s^*))}\right| = \left|\frac{h''(s^*)}{2(1 - \lambda_{p,x})}\right|.
\end{equation}
\end{theorem}

\begin{proof}
This is a standard result in fixed-point acceleration theory~\cite{stoer}. If $h$ has a fixed point $s^*$ with $h'(s^*) = \lambda \neq 1$ and $h$ is twice continuously differentiable near $s^*$, then the Steffensen map $T$ defined by~\eqref{eq:steffensen} satisfies $T(s^*) = s^*$ and $T'(s^*) = 0$, yielding quadratic convergence with the constant~\eqref{eq:K_steffensen}.
\end{proof}

\subsection{Numerical demonstration}

\begin{table}[h]
\centering
\begin{tabular}{cccc}
\toprule
Step $n$ & $s_n$ (Steffensen) & $|v_n - \cbrt{2}|$ & $|s_n - s^*|$ \\
\midrule
0 & $0.875000000000000$ & $2.71 \times 10^{-1}$ & $8.13 \times 10^{-2}$ \\
1 & $0.793949257650704$ & $7.90 \times 10^{-4}$ & $2.49 \times 10^{-4}$ \\
2 & $0.793700528604164$ & $8.32 \times 10^{-9}$ & $2.62 \times 10^{-9}$ \\
3 & $0.793700525984100$ & $2.22 \times 10^{-16}$ & $< 10^{-16}$ \\
\bottomrule
\end{tabular}
\caption{Steffensen-accelerated Pandrosion for $p=3$, $x=2$. Machine precision is reached in \textbf{3~steps} (each step requires 2~evaluations of $h$, i.e.\ 6 evaluations of $S_p$ total).}
\end{table}

\subsection{Comparison of quadratic constants}

Both Newton's method and Steffensen--Pandrosion converge quadratically, but with different constants.

\begin{table}[h]
\centering
\renewcommand{\arraystretch}{1.3}
\begin{tabular}{lcc}
\toprule
& \textbf{Newton} & \textbf{Steffensen--Pandrosion} \\
\midrule
Order & $2$ & $2$ \\
Constant ($p\!=\!3, x\!=\!2$) & $K_N = \dfrac{p-1}{2\alpha} \approx 0.794$ & $K_S \approx 0.013$ \\
Derivative required & Yes ($f'$) & No \\
Geometric origin & No & Yes \\
Steps to $10^{-15}$ & $5$ & $3$ \\
Evaluations of $S_p$ & --- & $6$ (total) \\
\bottomrule
\end{tabular}
\caption{Comparison of the two quadratically convergent methods for $\cbrt{2}$. Steffensen--Pandrosion has a much smaller quadratic constant ($K_S \ll K_N$), compensating for the double evaluation cost per step.}
\end{table}

\begin{remark}[A bridge between ancient and modern]
The Steffensen--Pandrosion method is remarkable: it takes a purely geometric, ancient construction (Pandrosion, order~1) and, through a universal algebraic acceleration (Aitken, 1926), produces a derivative-free, quadratically convergent method. This provides a concrete bridge between ancient geometric approximation and modern numerical analysis, demonstrating that the ``slow'' linear convergence of Pandrosion is not an intrinsic limitation but rather a platform for acceleration.
\end{remark}

%=====================================================================
\section{Scaling optimization: exploiting monotonicity}
\label{sec:scaling}
%=====================================================================

Theorem~\ref{thm:monotonicity} establishes that $\lambda_{p,x}$ is increasing in $x$: the larger the target $x$, the slower the convergence. For $x \gg 1$, the contraction ratio $\lambda_{p,x} \to 1^-$ (Proposition~\ref{prop:limits}), making direct application of Pandrosion's iteration impractical. We now show how to \emph{exploit} this very monotonicity to overcome it.

\subsection{The scaling principle}

The key observation is elementary: for any reference value $A > 1$ whose $p$-th root is known (or cheaply computable), we may write
\begin{equation}
  \label{eq:scaling}
  x^{1/p} = A^{1/p} \cdot \left(\frac{x}{A}\right)^{\!1/p}.
\end{equation}
Setting $x' = x/A$, Pandrosion's method is now applied to the \emph{reduced ratio} $x'$ instead of $x$. By choosing $A$ as the largest $p$-th power not exceeding $x$, we ensure $x' \in [1, (1+1/A^{1/p})^p)$, which is close to~$1$.

\begin{proposition}[Scaling reduction]
\label{prop:scaling}
Let $A = \lfloor x^{1/p} \rfloor^p$, so that $A^{1/p}$ is an integer and $x' = x/A \in [1, (1+1/\lfloor x^{1/p}\rfloor)^p)$. Then:
\begin{enumerate}
  \item The contraction ratio satisfies $\lambda_{p,x'} \ll \lambda_{p,x}$.
  \item The number of iterations drops from $O(\ln(\eps_0/\delta) / \ln(1/\lambda_{p,x}))$ to $O(\ln(\eps'_0/\delta) / \ln(1/\lambda_{p,x'}))$.
  \item The final result is recovered as $x^{1/p} = \lfloor x^{1/p} \rfloor \cdot (x')^{1/p}$.
\end{enumerate}
\end{proposition}

\begin{proof}
Part~(1) follows directly from Theorem~\ref{thm:monotonicity}: since $x' < x$ and $\lambda_{p,\cdot}$ is increasing, $\lambda_{p,x'} < \lambda_{p,x}$. Moreover, as $x \to \infty$, $x' \to 1^+$ so $\lambda_{p,x'} \to 0$ by Proposition~\ref{prop:limits}(1), while $\lambda_{p,x} \to 1^-$. Part~(2) follows from Proposition~\ref{prop:cout} applied to $x'$. Part~(3) is the factorization~\eqref{eq:scaling}.
\end{proof}

\begin{table}[h]
\centering
\begin{tabular}{ccccccc}
\toprule
$x$ & $\lambda_{3,x}$ & $A$ & $x' = x/A$ & $\lambda_{3,x'}$ & Reduction & Iterations \\
\midrule
$10$ & $0.615$ & $8$ & $1.25$ & $0.073$ & $8\times$ & $49 \to 8$ \\
$100$ & $0.890$ & $64$ & $1.5625$ & $0.145$ & $6\times$ & $213 \to 11$\\
$1000$ & $0.973$ & $729$ & $1.372$ & $0.103$ & $9\times$ & $>500 \to 9$ \\
$10000$ & $0.994$ & $9261$ & $1.080$ & $0.025$ & $39\times$ & $>500 \to 5$ \\
\bottomrule
\end{tabular}
\caption{Effect of scaling for $p = 3$ with $A = \lfloor x^{1/3} \rfloor^3$. The contraction ratio drops by up to $39\times$ and iteration counts collapse from hundreds to single digits. Iteration counts are for precision $10^{-10}$ with optimal starting point.}
\end{table}

\begin{remark}
The scaling reduction becomes more dramatic as $x$ grows: for $x = 10\,000$, the direct method requires over $500$ iterations while the scaled method needs only $5$. This is because $x' \to 1$ as $x \to \infty$, pushing $\lambda_{p,x'}$ toward the regime $\lambda \sim (p-1)(\alpha'-1)/2 \ll 1$ of Proposition~\ref{prop:limits}(1).
\end{remark}

\subsection{Geometric interpretation: Thales preconditioning}
\label{subsec:thales_preconditioning}

The scaling optimization admits an elegant geometric interpretation on Pandrosion's original Thales diagram.

\medskip

In the direct construction, one seeks to insert $p-1$ proportional means in a segment of ratio $x : 1$. When $x \gg 1$, this amounts to constructing parallels in a very elongated rectangle, where successive intercepts move slowly toward the target --- geometrically, the parallels are nearly horizontal.

\smallskip

The scaling strategy performs an \emph{homothety} on the construction: instead of starting from the unit segment, one starts from a segment of length $A^{1/p}$, chosen as the largest integer whose $p$-th power does not exceed $x$. The remaining construction now fills a much smaller gap: the rectangle has ratio $x' = x/A$ close to $1$, and the parallels converge rapidly.

\begin{proposition}[Double preconditioning]
\label{prop:preconditioning}
The scaling method combines two distinct acceleration mechanisms:
\begin{enumerate}
  \item \textbf{Initial residual reduction}: the reduced starting point $s_0^{\mathrm{opt}} = h_{x'}(1) = 1 - (x'-1)/(x'p)$ satisfies $|s_0^{\mathrm{opt}} - s'^*| \ll |s_0^{\mathrm{opt}} - s^*|$ since $x' \approx 1$ implies $s'^* = (x')^{-1/p} \approx 1$.
  \item \textbf{Intrinsic convergence acceleration}: each iteration contracts by $\lambda_{p,x'} \ll \lambda_{p,x}$, compounding the advantage over all subsequent steps.
\end{enumerate}
The total iteration count reduction is therefore superlinear in the ratio $\lambda_{p,x}/\lambda_{p,x'}$.
\end{proposition}

\begin{proof}
For mechanism~(1): $\eps'_0 = |v'_0 - \alpha'| \leq C \cdot |x' - 1| \to 0$ as $x' \to 1$, compared to $\eps_0 = |v_0 - \alpha| \geq C' \cdot |x - 1| \gg 1$ for the unscaled problem.

For mechanism~(2): the iteration count satisfies
\[
  n'_\delta \sim \frac{\ln(\eps'_0/\delta)}{\ln(1/\lambda_{p,x'})} \ll \frac{\ln(\eps_0/\delta)}{\ln(1/\lambda_{p,x})} \sim n_\delta,
\]
since both the numerator decreases (smaller $\eps'_0$) and the denominator increases (smaller $\lambda_{p,x'}$, hence larger $\ln(1/\lambda_{p,x'})$). These two effects multiply.
\end{proof}

\begin{remark}[From Pandrosion to a practical algorithm]
With scaling and Steffensen acceleration combined, Pandrosion's ancient geometric construction becomes a competitive root-finding algorithm: for any $x > 0$ and $p \geq 2$, one (i) reduces to $x' \approx 1$ via integer arithmetic, (ii) applies Steffensen--Pandrosion on $x'$ (converging in $2$--$3$ steps), and (iii) multiplies by $\lfloor x^{1/p} \rfloor$. The total cost is $4$--$6$ evaluations of $S_p$ plus one integer $p$-th root --- a purely algebraic, derivative-free computation whose origins lie in a ruler-and-parallels construction from fourth-century Alexandria.
\end{remark}

%=====================================================================
\section{Conclusion}
\label{sec:conclusion}
%=====================================================================

Starting from Pandrosion's geometric construction for $\cbrt{2}$, we have uncovered a universal structure: the iteration
\[
  s_{n+1} = 1 - \frac{x-1}{x \cdot S_p(s_n)},
\]
whose fixed point is $s^* = x^{-1/p}$ and whose output $v_n = x \cdot s_n^{p-1}$ converges to $x^{1/p}$. The convergence is linear, with exact contraction ratio
\[
  \lambda_{p,x} = \frac{(\alpha-1)\displaystyle\sum_{k=1}^{p-1} k\,\alpha^{p-1-k}}{\displaystyle\sum_{k=0}^{p-1}\alpha^k}, \qquad \alpha = x^{1/p},
\]
which depends only on the target $\alpha$ and the order $p$. This is the \textbf{convergence signature} of Pandrosion's method.

\medskip

Beyond this central result, we have established:
\begin{itemize}
  \item \textbf{Non-asymptotic bounds} (Theorem~\ref{thm:nonasymp}): $|s_n - s^*| \leq \Lambda^n |s_0 - s^*|$ for all $n$, with explicit $\Lambda$.
  \item \textbf{Double monotonicity} (Theorem~\ref{thm:monotonicity}): $\lambda_{p,x}$ is increasing in both $p$ and $x$.
  \item \textbf{Extension to $x < 1$} (Proposition~\ref{prop:x_small}): the method converges with \emph{oscillatory} contraction ratio $\lambda < 0$, with $|\lambda| < 1$.
  \item \textbf{Optimal starting point} (Proposition~\ref{prop:optimal_start}): $s_0 = h(1) = 1 - (x-1)/(xp)$ minimizes the initial residual.
  \item \textbf{Steffensen acceleration} (Theorem~\ref{thm:steffensen}): Aitken's $\Delta^2$ process transforms Pandrosion's linear method into a quadratically convergent, derivative-free method with constant $K_S \ll K_{\text{Newton}}$.
  \item \textbf{Scaling optimization} (Proposition~\ref{prop:scaling}): exploiting the monotonicity of $\lambda_{p,x}$, a simple change of variable reduces large-$x$ problems to $x' \approx 1$, cutting the contraction ratio by up to $39\times$ and collapsing iteration counts from hundreds to single digits.
\end{itemize}

\medskip

The coexistence of \emph{five} distinct residual profiles (Pandrosion, bisection, secant, Newton, Steffensen--Pandrosion) for the same target $\alpha = \cbrt{2}$ illustrates the central thesis: \textbf{the residual profile is a property of the method, not of the number.}

\medskip

\textbf{Perspectives.} It would be natural to study generalized nonlinear Pandrosion constructions (with other polynomials replacing $S_p$), to investigate whether higher-order Steffensen variants produce even faster convergence, or to analyze the topological structure of the basin of attraction for $x < 1$ in higher dimensions. The close relationship between Pandrosion's iteration and the arithmetic of roots also suggests connections with continued fraction algorithms that deserve further exploration.

%=====================================================================
\begin{thebibliography}{9}

\bibitem{mactutor_pandrosion}
J.~J. O'Connor and E.~F. Robertson,
``Pandrosion of Alexandria,''
\textit{MacTutor History of Mathematics}, University of St~Andrews, 2020.
\url{https://mathshistory.st-andrews.ac.uk/Biographies/Pandrosion/}

\bibitem{mactutor_cube}
J.~J. O'Connor and E.~F. Robertson,
``Doubling the cube,''
\textit{MacTutor History of Mathematics}, University of St~Andrews, 2002.
\url{https://mathshistory.st-andrews.ac.uk/HistTopics/Doubling_the_cube/}

\bibitem{pappus}
Pappus of Alexandria,
\textit{Collectio}, Book~III (circa 340~AD).
Critical edition: F.~Hultsch, \textit{Pappi Alexandrini Collectionis quae supersunt},
Weidmann, Berlin, 1876--1878.

\bibitem{wantzel}
P.~Wantzel,
``Recherches sur les moyens de reconna\^itre si un probl\`eme de G\'eom\'etrie peut se r\'esoudre avec la r\`egle et le compas,''
\textit{Journal de Math\'ematiques Pures et Appliqu\'ees}, s\'er.~1, vol.~2, pp.~366--372, 1837.

\bibitem{stoer}
J.~Stoer and R.~Bulirsch,
\textit{Introduction to Numerical Analysis},
3rd~ed., Springer, New York, 2002.
Chap.~5: fixed-point iterations, convergence rates, and Aitken acceleration.

\bibitem{numrecipes}
W.~H. Press, S.~A. Teukolsky, W.~T. Vetterling, and B.~P. Flannery,
\textit{Numerical Recipes: The Art of Scientific Computing},
3rd~ed., Cambridge University Press, 2007.

\bibitem{steffensen}
J.~F. Steffensen,
``Remarks on iteration,''
\textit{Skandinavisk Aktuarietidskrift}, vol.~16, pp.~64--72, 1933.

\bibitem{aitken}
A.~C. Aitken,
``On Bernoulli's numerical solution of algebraic equations,''
\textit{Proceedings of the Royal Society of Edinburgh}, vol.~46, pp.~289--305, 1926.

\end{thebibliography}

\end{document}
